
In a regression setting, under a squared error loss, we estimate $\E[Y \given \Xvec = \xvec]$, for example, through:
\begin{enumerate}
    \item \emph{Kernel-based methods}:
          \begin{equation}
              \widehat{\E[Y \given \Xvec = \xvec]} = \Avg(y_i \given \xvec_i \in N_k(\xvec)).
          \end{equation}
    \item \emph{Linear regression}:
          \begin{equation}
              \widehat{\E[Y \given \Xvec = \xvec]} = \hat{\beta}_0 + \hat{\beta}_1 x_1 + \cdots + \hat{\beta}_p x_p = \xvec^\t \hat{\betavec}.
          \end{equation}
\end{enumerate}

In a classification setting, under a $0$-$1$ loss, we estimate $\P(1 \given \Xvec = \xvec)$, for example, through:
\begin{enumerate}
    \item \emph{$k$-nearest neighbors classifier} (non-parametric).
    \item \emph{Logistic regression} (parametric).
\end{enumerate}
Those methods will provide an estimate of $\P(1 \given \Xvec = \xvec)$, namely $\widehat{\P(1 \given \Xvec = \xvec)}$, or more generally:
\begin{equation}
    \P(Y = j \given \Xvec = \xvec), \quad j = 1, \ldots, C,
\end{equation}
where $C$ is the number of classes.

\section{\texorpdfstring{$k$-nearest neighbors classifier}{k-nearest neighbors classifier}}

Let $k$ be a positive integer and $\tuple{\xvec_i, y_i}, i = 1, \ldots, n$ be training data. Then, $\P(Y = j \given \Xvec = \xvec_0)$, with $\xvec_0$ as a realization of $\Xvec$, is estimated as follows:
\begin{enumerate}
    \item Identify the $k$ training data that are closest to $\xvec_0$. Denote these as $N_k(\xvec_0)$.
    \item Estimate the following quantity:
          \begin{equation}
              \widehat{\P(Y = j \given \Xvec = \xvec_0)} = \frac{1}{k} \sum_{i \in N_k(\xvec_0)}\hspace{-8pt}\oone(y_i = j), \quad j = 1, \ldots, C,
          \end{equation}
          where $\oone(\blank)$ is the indicator function.
    \item Classify $\xvec_0$ to the class with the largest estimated probability ($0$-$1$ loss, threshold $0.5$ for binary classification).
\end{enumerate}

The choice of $k$ is crucial for the performance of the method. This is connected to the bias-variance trade-off. We identify two cases:
\begin{itemize}
    \item With \emph{large values of $k$} we have a simpler decision surface, thus less flexible, but it may generalize well to new data. However, we also have high bias and low variance.
    \item With \emph{small values of $k$} the surface is more local, more flexible/jiggly and may not generalize well to new data. The bias-variance relationship is reversed, with low bias and high variance.
\end{itemize}

\subsection{Error rate on test data}

In order to find the optimal trade-off between bias and variance, the simplest approach is to calculate the error rate on unseen test data, \ie:
\begin{enumerate}
    \item Calculate probabilities on test data for a range of values of $k$.
    \item Use the response values of the test data to calculate error rate.
    \item Choose the value of $k$ that minimizes the error rate.
\end{enumerate}

\section{Logistic regression}

This method is the equivalent of linear regression for classification problems. Why can't we use linear regression for classification problems? Let's discuss some examples to understand why.

\begin{examples}\leavevmode
    \begin{enumerate}
        \item Suppose we want to predict medical condition of patients arriving at emergency based on their symptoms. Assume that the response $Y$ can take only three possible values:
              \begin{equation}
                  Y = \begin{cases}
                      1 & \text{if the patient has a stroke},            \\
                      2 & \text{if the patient has a drug overdose},     \\
                      3 & \text{if the patient has an epileptic seizure}
                  \end{cases}
              \end{equation}
              and that the symptoms are represented by a vector $\xvec$. If we use linear regression to predict $Y$ based on $\xvec$, \ie, we try to fit a line to the data, we have two problems:
              \begin{itemize}
                  \item We are implying an ordering between the values of the response
                  \item We are implying an equal difference between pairs of response values, \eg, the difference between $1$ and $2$ is the same as the difference between $2$ and $3$.
              \end{itemize}
        \item Let us now consider only a binary response:
              \begin{equation}
                  Y = \begin{cases}
                      1 & \text{if the patient has a drug overdose}, \\
                      0 & \text{if the patient has a stroke}.
                  \end{cases}
              \end{equation}
              Since $Y$ is binary, it follows a Bernoulli distribution:
              \begin{equation}
                  f(y) = \begin{cases}
                      p     & \text{if } y = 1, \\
                      1 - p & \text{if } y = 0.
                  \end{cases}
              \end{equation}
              Therefore:
              \begin{equation}
                  \E[Y \given \Xvec = \xvec] = 1 \cdot \P(1 \given \xvec) + 0 \cdot \P(0 \given \xvec) = \P(1 \given \xvec).
              \end{equation}
              This means that linear regression makes the assumption:
              \begin{equation}
                  \E[Y \given \Xvec = \xvec] = \underbracket{\P(1 \given \xvec)}_{\in [0,1]} = \beta_0 + \beta_1 x_1 + \cdots + \beta_p x_p = \underbracket{\xvec^\t \betavec}_{\mathclap{\in (-\infty, +\infty)}},
              \end{equation}
              which is not possible since the left-hand side is a probability and thus must be between $0$ and $1$.
    \end{enumerate}
\end{examples}

For logistic regression, we assume:
\begin{equation}
    \P(1 \given \xvec) = \frac{e^{\xvec^\t\betavec}}{1 + e^{\xvec^\t \betavec}} \in (0,1).
\end{equation}
This function is called \emph{logistic function} or \emph{sigmoid}.

\begin{figure}[tb]
    \centering
    \begin{tikzpicture}
        \begin{axis}[
                axis lines=middle,
                enlargelimits,
                xlabel=$\xvec^\t \betavec$,
                ylabel=$\P(1 \given \xvec)$,
                xtick={0},
                ytick={0, 0.5, 1},
                width=.8\textwidth,
                height=.4\textwidth,
                legend style={
                        at={(.29\textwidth, 5.5cm)},
                        anchor=north
                    },
                legend cell align=left,
                smooth
            ]
            \addplot[domain=-10:10, samples=100, thick] {exp(x) / (1 + exp(x))};
            \addplot[domain=-10:10, samples=100, thick, dashed] {exp(-x) / (1 + exp(-x))};
            \legend{\text{$\beta_0 = 0, \beta_1 = 1$}, \text{$\beta_0 = 0, \beta_1 = -1$}}
        \end{axis}
    \end{tikzpicture}
    \caption{Logistic or \emph{sigmoid} function}
\end{figure}

\begin{example}
    For $\betavec = \tuple{\beta_0, \beta_1}$ the explicit function is:
    \begin{equation}
        \P(1 \given \xvec) = \frac{e^{\beta_0 + \beta_1 x_1}}{1 + e^{\beta_0 + \beta_1 x_1}}.
    \end{equation}
\end{example}

We can interpret the coefficients in terms of \emph{odds}:
\begin{equation}
    \frac{\P(1\given \xvec)}{\P(0\given \xvec)} = \frac{\P(1\given \xvec)}{1 - \P(1\given \xvec)} = \frac{\frac{e^{\xvec^\t\betavec}}{1 + e^{\xvec^\t\betavec}}}{1 - \frac{e^{\xvec^\t\betavec}}{1 + e^{\xvec^\t\betavec}}} = e^{\xvec^\t\betavec}\negthinspace.
\end{equation}
If we take the $\log$ we get the \emph{log-odds} (or \emph{logit}):
\begin{equation}
    \log\left(\frac{\P(1\given \xvec)}{1 - \P(1\given \xvec)}\right) = \xvec^\t\betavec.
\end{equation}

We have now reduced the problem to the estimation of $\betavec$ (parameter estimation). To do so, we rely on the method of \emph{maximum likelihood estimation} (\acs{mle}). The likelihood function corresponds to the likelihood of the specific configuration of data we have observed:
\begin{equation}
    \L(\betavec) = \prod_{i = 1}^n \P(1 \given \xvec_i)^{y_i} (1 - \P(1 \given \xvec_i))^{1 - y_i}\negthinspace.
\end{equation}
We can also compute the \emph{log-likelihood}:
\begin{equation}
    \begin{aligned}
        \l(\betavec) = \log\L(\betavec) & = \sum_{i = 1}^n y_i\log \P(1 \given \xvec_i) + (1 - y_i)\log\bigl(1 - \P(1 \given \xvec_i)\bigr) \\
                                     & = \sum_{i = 1}^n \left(y_i\log\left(\frac{\P(1 \given \xvec_i)}{1 - \P(1 \given \xvec_i)}\right) + \log\bigl(1 - \P(1 \given \xvec_i)\bigr)\right)\\
                                     & = \sum_{i = 1}^n \left(y_i \xvec_i^\t \betavec - \log\bigl(1 + e^{\betavec^\t \xvec_i}\bigr)\right).
    \end{aligned}
\end{equation}

To find the $\betavec$ that maximizes the log-likelihood, \ie, solving:
\begin{equation}
    \pdv{\l(\betavec)}{\beta_j} = 0, \quad j = 0, \ldots, p,
\end{equation}
there is no closed-form solution, unlike for linear regression, for which we have:
\begin{equation}
    \betavech = (X^\t X)^{-1} X^\t Y,
\end{equation}
so the equations can only be solved numerically (for example, using an adaptation of the Newton-Raphson method). 

At the end of the estimation we get different quantities of interest:
\begin{enumerate}
    \item The estimated coefficients of the model: $\betavech_0, \betavech_1, \ldots, \betavech_p$.
    \item The $p$-values and $H_0$ hypothesis tests for the significance of each coefficient:
    \begin{equation}
        H_0 = \begin{cases}
            \beta_j = 0,\\
            \beta_j \neq 0,
        \end{cases}\quad \text{for } j = 0, \ldots, p.
    \end{equation}
    \item The prediction:
    \begin{equation}
        \widehat{\P(1 \given \xvec)} = \frac{e^{\betavech^\t \xvec}}{1 + e^{\betavech^\t \xvec}}.
    \end{equation}
    \item Classification, \acs{roc} curve on test data, etc.
\end{enumerate}